\documentclass[11pt,a4paper]{article}
\usepackage[margin=24mm]{geometry}
\usepackage[T1]{fontenc}
\usepackage{lmodern,amsmath,amssymb,amsthm,mathtools,booktabs,array,microtype,enumitem,fancyhdr,graphicx}
\usepackage[dvipsnames]{xcolor}
\usepackage[colorlinks=true,linkcolor=MidnightBlue,urlcolor=MidnightBlue,citecolor=MidnightBlue]{hyperref}
\usepackage{bookmark}
\hypersetup{pdftitle={Bounded-component acquisition of an unknown many-body quantum register},pdfauthor={Astra P1}}
\setlength{\parskip}{5pt}\setlength{\parindent}{0pt}
\setlist{itemsep=3pt,topsep=4pt}
\pagestyle{fancy}\fancyhf{}
\fancyhead[L]{\small Astra P1 | Bounded-component quantum acquisition}
\fancyhead[R]{\small 5 November 2026}\fancyfoot[C]{\thepage}
\setlength{\headheight}{14pt}
\newtheorem{theorem}{Theorem}\newtheorem{lemma}[theorem]{Lemma}\newtheorem{corollary}[theorem]{Corollary}
\newcommand{\id}{\operatorname{id}}\newcommand{\Tr}{\operatorname{Tr}}\newcommand{\osc}{\operatorname{osc}}
\newcommand{\norm}[1]{\left\lVert #1\right\rVert}\newcommand{\ket}[1]{\lvert #1\rangle}\newcommand{\bra}[1]{\langle #1\rvert}
\newcommand{\cA}{\mathcal A}\newcommand{\cC}{\mathcal C}\newcommand{\cD}{\mathcal D}\newcommand{\cR}{\mathcal R}
\newcommand{\ee}{\mathrm e}
\begin{document}
\hypersetup{pageanchor=false}
\begin{titlepage}
\vspace*{12mm}{\Large\color{MidnightBlue} ASTRA P1 SCIENTIFIC RETURN\par}
\vspace{10mm}{\huge\bfseries Bounded-component acquisition\newline of an unknown many-body\newline quantum register\par}
\vspace{7mm}{\Large Continuous fields, two-level receivers,\newline fixed local response, and an explicit decoder\par}
\vspace{9mm}{\large 5 November 2026, Australia/Brisbane\par}
\vspace{11mm}
\textbf{Completed operation.} Transfer the initial, arbitrary, unencoded quantum register of a known interacting source into receiver qubits, including its correlations with every external reference. The receiver uses only the actual coherently delivered output ports. Every material capture cell is two-dimensional; its coupling norm and control-response poles stay fixed as register size and target accuracy change.

\textbf{Price.} Hardware, propagation distance, preparation, calibration precision, memory lifetime, and quantum processing grow. Explicit conservative bounds are polynomial, with very large exponents. This is a construction in a declared calibrated chiral Markov interface, not a demonstrated device or a competitive engineering estimate.

\textbf{Boundary earned with the construction.} Fixed local linewidth is compatible with increasing collective acceptance. Fixed inaccessible collection loss or a fixed common spectral blind subspace is a different resource: even the noninteracting Weyl source has an exponential full-register Bell-fidelity ceiling.

\vfill The preceding continuous-Hamiltonian theorem, its finite collector, source-plus-field comparison, reverse-reset decoder, stopping certificate, diagonal extension, and nested predecessor package are preserved byte-for-byte in the accompanying archive.
\end{titlepage}
\hypersetup{pageanchor=true}
\begin{abstract}
The previous scalable acquisition theorem required local collectors whose peak strength and cycle deadline became increasingly demanding. This report replaces those collectors by a continuous passive receiver and a finite qubit compiler. A static splitter dilutes each physical output among polynomially many rails. Each rail contains identically controlled two-level absorbers. On the sector with at most one photon per rail, this physical receiver agrees exactly with a passive linear network. A uniform factorial-moment bound controls the remaining sector without postselection. The linear network stores explicitly derived warped Laguerre modes; a quantitative projection bound gives polynomial approximation of the temporal modes needed by a reset/export instrument. The stored qubits become binary occupation registers directly, and bounded local gates perform spatial recombination, temporal-mode conversion, local Kraus conversion, and the existing reverse-reset decoder.

All control coefficients may be produced by fixed three-pole actuators. A constructive parameter recipe resolves the dependence of switching guards on photon cutoff and receiver size; exact instantaneous switch-off is unnecessary. The operation is deterministic and reference-uniform. Its ideal theorem still requires actual pure coherent ports, lossless transport, and sufficiently accurate receiver operation. For the same Weyl source class, a complementary-environment argument and an exactly tensorizing recovery semidefinite program prove a nonzero loss floor and exponential many-input obstruction. A finite-time spectral version establishes a corresponding boundary for a fixed aggregate inaccessible frequency band. Neither converse is a zero-capacity theorem.
\end{abstract}
\setcounter{tocdepth}{1}
{\small\setlength{\parskip}{0pt}\tableofcontents}\newpage

\section{Decision, consequence, and ordinary comparison}

The destination is a quantum interface between an interacting experiment and a later quantum computation. It must transfer the existing unknown register, not learn a classical model of it, wait for a fresh payload, or infer a decoder from source forgetting alone. Once the acquired channel is within diamond distance $\varepsilon$ of the identity, composition with \emph{any} later quantum experiment increases that experiment's channel error by at most $\varepsilon$. This is the general consequence earned here.

Fine temporal bins are synthesized \emph{after} coherent storage. They no longer impose physical cycling deadlines. Continuous capture uses a fixed family of local responses, while numerical resolution is paid by more stationary coordinates and a longer ordinary quantum circuit.

Passive linear quantum memories, matched temporal modes, writing/storage separation, and optical unitary decomposition are established machinery \cite{yamamoto,nurdin,reck}. An equally informed ordinary architecture with ideal oscillator memories can use the same modal representation and decoder without the rail overhead. We do not claim a better oscillator memory. The stronger resource statement here is a compiler to two-level capture cells with bounded material-system coupling norms, an independently initialized finite-front interface, and a complete arbitrary-input acquisition circuit. A source-side parallel SWAP or a sender-shaped photon can be substantially simpler when the additional source control is granted; it is absent from this task.

Fixed local response means bounded coupling coefficients, fixed linewidth, and fixed actuator poles. It does not mean compact Fourier support, a fixed total accepted optical band, or a bound on the microscopic white-field interaction Hamiltonian. The governing Astra N1732 card keeps the same distinction \cite{astra1732}. No microscopic ultraviolet completion or fixed-noise front-end threshold is asserted.

\section{Source, receiver, and operational promise}
\label{sec:contract}

Let $S=(\mathbb C^q)^{\otimes n}$ with fixed $q\ge2$. The input $\rho_{SR}$ is arbitrary, unencoded, and possibly entangled with an inaccessible reference $R$. The receiver is initially independent. There is no source-side gate, feedback, protected initial encoding, or later access to $S$.

The source has known static $k$-local Hamiltonian $H=\sum_{X\in\mathcal E}h_X$, with fixed $k$, polynomially many terms $m=|\mathcal E|$, supplied matrices, and charged precision. Its actual independent vacuum ports are
\begin{equation}
 L_{i\mu}=\frac{\sqrt r}{q}W_{\mu,i},\qquad \mu=1,\ldots,q^2-1,
\end{equation}
where $W_0=I,W_1,\ldots,W_{q^2-1}$ is an orthogonal Weyl unitary basis. Every modeled output reaches the receiver coherently with known phase and label. The ideal theorem has no additional inaccessible source noise. Its reduced equation is
\begin{equation}
 \mathcal L(\rho)=-i[H,\rho]+r\sum_i(\mathcal R_i-\id)(\rho),\qquad
 \mathcal R_i(\rho)=\frac{I_i}{q}\otimes\Tr_i\rho.
\end{equation}
This equation alone does not supply the ports. Define
\begin{align}
 F&=n(q^2-1),&\nu&=r/q^2,&c&=r(1-q^{-2}),&R_p&=F\nu=nc,\\
 \beta&=\max_i\sum_{X\ni i}(|X|-1)\osc(h_X),&&
 \chi=\max_i\sum_{X\ni i}\osc(h_X),&&\kappa=r-\beta.
\end{align}
The physical count rate $c$ differs from synthetic reset rate $r$; $R_p$ denotes total photon rate, not the reference system. On-site fields enter $\chi$ and the execution bill but not $\beta$.

Hardware consists of calibrated one-way dispersionless paths, static directional couplers, two-level capture and storage cells, isolated cold reset ports, and bounded local one- and two-qubit Hamiltonians. Capture uses $L_j(t)=g_j(t)\sigma_j^-$, with $|g_j|\le g_0$; local gate norm is at most a fixed $\Omega_0$; actuator poles are fixed at $B>0$. These constants do not grow with $n$ or $\varepsilon^{-1}$. Extra cells, control lines, cold-bath capacity, and parallel power are charged hardware. Uncontrolled errors are priced in Section~\ref{sec:errors}.

The desired deterministic channel $\cA:S\to S'$ satisfies
\begin{equation}
 \norm{\cA-\id}_\diamond\le\varepsilon.
 \label{eq:target}
\end{equation}
Trace and diamond norms are unhalved. Bell fidelity is overlap with the normalized target Bell state. Equation~\eqref{eq:target} preserves every input/reference correlation and has success probability one.

\begin{theorem}[Fixed local response, growing finite hardware]
\label{thm:positive}
Under this interface, if $\kappa>0$, an explicitly computable causal receiver satisfies \eqref{eq:target}. Capture cells have dimension two and coupling norm at most $g_0$; local gates have fixed norm and response-pole bounds. Hardware, numerical preprocessing, program length, preparation, propagation delay, storage, and execution time are polynomial in the supplied size, norm, inverse-margin, inverse-rate and inverse-accuracy parameters, for fixed $q,k$ and fixed hardware constants. Receiver error and calibration budgets must be supplied. The operation assumes neither an ideal field buffer nor a field-to-qubit transducer.
\end{theorem}

The conservative polynomial can be enormous. This theorem does not establish useful economics, fixed uncontrolled capture-loss tolerance, or acquisition from every dilation of the same reduced generator. The condition $r>\beta$ remains sufficient, not necessary.

A directly relevant family has bounded interaction degree and strength, fixed $r$, and a fixed positive margin $\kappa$. Then every local source scale as well as every receiver strength and response pole stays bounded as the register grows. No growing source export rate is used to obtain this family.

\section{Uniform field statistics and exact qubit compilation}

A resolved photon history has source Kraus operator equal to a scalar times a product of Weyl unitaries and Hamiltonian propagators. Thus $K_\gamma^\dagger K_\gamma$ is the scalar history probability times identity. Every original port has an input-independent Poisson count of rate $\nu$, uniformly with arbitrary reference and interacting $H$.

Split each original port into $P$ equal rails, with $P$ a power of two and all other inputs independent vacuum. In first quantization, an $s$-photon input acquires spatial factor $u^{\otimes s}$, where $u=P^{-1/2}\sum_p\ket p$, regardless of temporal entanglement. Each pair shares a rail with probability $1/P$. For active source-time window $A$,
\begin{equation}
 p_{\rm rail}\le\frac{F\nu^2A^2}{2P}.
 \label{eq:rail}
\end{equation}
This is an operator-uniform bound. The good/bad projection is only a proof device; no count or collision test is performed.

Each rail has $N$ initially ground-state absorbers. If at most one photon enters that entire rail over the active window, excitation conservation confines its field and material systems to vacuum plus one excitation. There, the qubit and oscillator interactions agree exactly, including creation terms on reachable states, delays, readout, and reference entanglement. On the other sector the difference of two isometries has norm at most two. Therefore
\begin{equation}
 \norm{\cC_{\rm qubit}-\cC_{\rm linear}}_\diamond\le4\sqrt{p_{\rm rail}}.
 \label{eq:qubitcompiler}
\end{equation}
More generally, any emitted-field isometry $V$ satisfying $V^\dagger N_j(N_j-1)V\le f_jI$ admits the same compiler with error $4\sqrt{\sum_j f_j/(2P_j)}$. This permits arbitrary photon-number superpositions, temporal correlations, and inter-port entanglement.

The oscillator is a calculational model. Every actual material cell is two-dimensional, with $\norm{g\sigma^-}=|g|\le g_0$. No untruncated bosonic annihilator is assigned a bounded norm. The emitted field is never first placed in an assumed buffer: the actual chiral receiver executes these interactions chronologically.

\section{The physical field-to-register operation}
\label{sec:map}

\subsection{Static splitters and real paths}

Two co-propagating chiral fields of speed $v$ interact over a fixed region by
\begin{equation}
 H_{\rm path}=-iv\int\psi^\dagger\partial_x\psi\,dx
 +\int J(x)(\psi_1^\dagger\psi_2+\psi_2^\dagger\psi_1)\,dx,
 \qquad |J(x)|\le J_0.
\end{equation}
The characteristic solution is $\exp[-i\theta\sigma_x]$ and a real transit delay, with $\theta=\int J(x)dx/v$. Set $\theta=\pi/4$ for a 50:50 directional coupler. Its local coefficient and size stay fixed; a tree uses $P-1$ such components per original port.

Logical depth is not physical radius. With leaves at pitch $2a$, a safe planar tree has equal root-to-leaf length $a(P-1)+a\log_2P$, total tree wire $O(aP\log P)$, and $P$ length-$N$ absorber rails with wire $O(aPN)$. Per-port internal delay is $O(a(P+N)/v)$ and total wire is
\begin{equation}
 O\bigl(Fa[P\log P+PN]\bigr).
\end{equation}
External delivery and placement of the modules are additional specified paths. A side-by-side planar arrangement has polynomial area and distances; actual delivery delay and loss are charged.

Let $d_{pj}$ be the calibrated source-to-node delay. Node $(p,j)$ uses $g(t-d_{pj})$. Operations on distinct nodes and different passing slices commute, so the spacetime circuit can be reordered into a spatial cascade in retarded source time. This does not grant field reversal or zero-delay hardware. At the common final clock, stored registers coexist; residual travelling field is traced, not counted as free memory.

\subsection{Exact stored modes and finite startup}

On the linear comparison sector, set $A_0=-I/2-J_{<}$, where $J_{<}$ has ones strictly below the diagonal. Write
\begin{equation}
 \gamma(s)=g(s)^2,\quad u(s)=\int_0^s\gamma(v)dv,\quad
 U=u(t_2),\quad x(s)=U-u(s).
\end{equation}
The local input-output equations give
\begin{equation}
 a(t_2)=\ee^{A_0U}a(0)-\int_0^{t_2}g(s)\ee^{A_0x(s)}\mathbf1\,dB(s).
 \label{eq:capture}
\end{equation}
The identity $[\ee^{A_0x}\mathbf1]_j=\ee^{-x/2}L_{j-1}(x)$ gives the actual stored field kernels. This is the acquired transduction map.

Finite initial vacuum coordinates complete the orthogonal modes:
\begin{equation}
 \int_0^U\ee^{A_0x}\mathbf1\mathbf1^\dagger\ee^{A_0^\dagger x}dx
 =I-\ee^{A_0U}\ee^{A_0^\dagger U}.
 \label{eq:gramvac}
\end{equation}
Representing the missing interval $x>U$ by mathematical vacuum is only a representation of these finite initial coordinates. No infinitely long rising photon, empty output prehistory, or prior field buffer is supplied. Before source time zero the signal couplings vanish, so older source output passes through without entangling the receivers.

\subsection{Fixed-response controls and the leading edge}

Use three cascaded first-order filters with fixed pole $B$ and step response
\begin{equation}
 F_3(x)=1-\ee^{-x}(1+x+x^2/2)\ (x\ge0),\qquad F_3(x)=0\ (x<0).
\end{equation}
For a bounded command held until $t_1$, the actual coupling is
\begin{equation}
 g(s)=g_0\{F_3(Bs)-F_3(B(s-t_1))\}.
 \label{eq:actuator}
\end{equation}
It has fixed peak, slew, and response poles; numerical binning imposes no modulation at the bin rate. On the useful window $[t_0,t_1]$, with $0<t_0\le1/B$,
\begin{equation}
 \gamma(s)\ge\alpha:=\frac{g_0^2\ee^{-2}(Bt_0)^6}{36},\qquad
 |\gamma'(s)|\le K_1:=2g_0^2B.
 \label{eq:rise}
\end{equation}
After the turnoff command, $|g(t_1+w)|\le2g_0\ee^{-Bw/2}$ and
\begin{equation}
 \int_{t_1}^{\infty}|g(s)|^2ds\le\frac{33g_0^2}{16B}.
 \label{eq:tailarea}
\end{equation}

Let $t_1=t_0+T$. The source did not wait for startup. Tracing the early output yields
\begin{equation}
 \mathcal E_{[0,t_0]}=\ee^{-R_pt_0}\mathcal U_H(t_0)
 +(1-\ee^{-R_pt_0})\mathcal N
\end{equation}
for a channel $\mathcal N$, so $\norm{\mathcal E_{[0,t_0]}-\mathcal U_H(t_0)}_\diamond\le2R_pt_0$. Recover the state at $t_0$, then apply $U_H(t_0)^\dagger$ on receiver hardware. The target is the original input. Early and trailing photons stay in the capture calculation and are later discarded as unwanted modes; there is no abrupt physical temporal filter.

\section{Polynomial approximation of the required temporal modes}

Divide $[t_0,t_1]$ into $K$ bins of width $\Delta=T/K$. Let $e_j=\Delta^{-1/2}1_{\mathrm{bin}\ j}$ and let $E$ embed these orthonormal modes into the complete one-particle input space, including initial vacuum coordinates. Let $P_N$ project onto the stored coordinates. Calculate
\begin{equation}
 B_{\ell j}=\langle f_\ell,e_j\rangle,\quad G=B^\dagger B=E^\dagger P_NE.
\end{equation}
If $\norm{I-G}\le\zeta<1$, then $W=P_NEG^{-1/2}$ is an orthonormal family inside actual stored modes. Principal angles give a unitary $V$ with $VE=W$ and
\begin{equation}
 \norm{V-I}=\sqrt{2-2\sqrt{\lambda_{\min}(G)}}\le2\sqrt\zeta.
\end{equation}
This $V$ is only an analysis comparison, not a physical operation on unavailable fields.

For an explicit size bound, put $X_0=g_0^2[t_1+33/(16B)]$. A bin becomes $h(x)=e_j(t(x))/\sqrt{\gamma(t(x))}$. Smooth its edges linearly over physical width $w\le\Delta/2$. The transformed smoothed function obeys
\begin{align}
 \norm{h-h_s}_2^2&\le2w/(3\Delta),\\
 \norm{h_s'}_2^2&\le4/(\alpha^2\Delta w)+K_1^2/(2\alpha^4).
\end{align}
Laguerre functions diagonalize $\mathsf A f=-(xf')'+xf/4$ with eigenvalues $j+1/2$. Its quadratic form gives
\begin{equation}
 \norm{(I-P_N)h}_2^2\le\frac{4w}{3\Delta}
 +\frac{8X_0/(\alpha^2\Delta w)+X_0K_1^2/\alpha^4+X_0/2}{N+1/2}.
\end{equation}
The sum of all $K$ tails bounds the positive Gram deficit. With $w=\Delta\zeta/(8K)$, it suffices to choose
\begin{equation}
 N\ge K,\qquad
 N\ge\frac{256X_0K^2}{\alpha^2\Delta^2\zeta^2}
       +\frac{4X_0KK_1^2}{\alpha^4\zeta}
       +\frac{2X_0K}{\zeta}.
 \label{eq:modecount}
\end{equation}
The startup coupling $\alpha$ is inverse-polynomial when $t_0$ is. Its large inverse powers are explicitly charged.

The actual total photon count over $A$ is Poisson with mean $R_pA$. For tail allowance $p_Q$, choose
\begin{equation}
 Q=\left\lceil(\ee-1)R_pA+\log(1/p_Q)\right\rceil.
 \label{eq:cap}
\end{equation}
The overflow probability is at most $p_Q$. On the $\le Q$ sector, $\norm{\Gamma(V)-I}\le2Q\sqrt\zeta$. Extracting stored $W$ rather than exact bin modes therefore costs at most $4Q\sqrt\zeta$, plus a conservative $4\sqrt{p_Q}$ cutoff allowance. The one-particle map is a direct sum across ports; $Q$ already counts all photons.

This comparison is made on the \emph{actual} source output before normalizing a later good-history model. It includes early and trailing photons and arbitrary references. This order avoids an uncharged conditional Poisson-tail factor.

Preprocessing is constructive. For $x\ge0$, $|\ee^{-x/2}L_j(x)|\le1$, equivalently a diagonal oscillator displacement matrix element. The derivative recurrence gives a bound $j+1/2$. With the explicit smooth controls, polynomial quadrature computes the overlaps to polynomial precision. The Gram inverse is well conditioned since $\lambda_{\min}(G)\ge1-\zeta$. Dense numerical matrices have dimension at most $N$, not $q^n$.

\section{Actual source-plus-field comparison and reset/export}

Pass the following comparisons through the already defined exact-bin extraction and decoder. The preserved full-history argument compares the actual source-plus-field isometry over $T$ with noise-only bins followed by exact Hamiltonian bin unitaries. Each jump crosses at most $\Delta$ of Hamiltonian evolution and contributes at most $\chi\Delta$. Orthogonal complete histories and the Poisson second moment give
\begin{equation}
 \eta_{\rm split}\le2\chi\Delta\sqrt{R_pT+(R_pT)^2}.
 \label{eq:split}
\end{equation}
This controls the coherent field and source, not only the reduced generator. No source pause or $R_p\Delta\le1$ assumption is used.

In the split model, analytically restrict to at most one physical photon per site/bin. Its input-independent bad probability is at most $nc^2T\Delta/2$. The normalized isometric comparison channel is within
\begin{equation}
 \eta_{\rm bin}\le2\sqrt{nc^2T\Delta/2}.
 \label{eq:onebin}
\end{equation}
There is no physical postselection. The comparison has flat temporal modes and local Kraus operators
\begin{equation}
 K_0=I/\sqrt{1+c\Delta},\qquad K_\mu=\sqrt\Delta L_\mu/\sqrt{1+c\Delta}.
\end{equation}
Its channel is exactly $(1-p)\id+p\mathcal R$, with
\begin{equation}
 p=\frac{r\Delta}{1+c\Delta},\qquad
 r_{\rm eff}=-\frac{\log(1-p)}{\Delta}\ge r-r^2\Delta/2.
 \label{eq:reff}
\end{equation}
Take $r\Delta\le1$ and $\Delta\le\kappa/r^2$, so $p\le1$ and $r_{\rm eff}-\beta\ge\kappa/2$.

With $c_0=1-p+p/q^2$, an explicit coherent environment isometry is
\begin{align}
 V_E\ket0&=\sqrt{(1-p)/c_0}\ket{\emptyset}
              +\frac{\sqrt p}{q\sqrt{c_0}}\ket{\Phi_q}_{AB},\\
 V_E\ket\mu&=q^{-1/2}\sum_{a,b}(W_\mu)_{ab}^*\ket{a,b},\quad\mu\ne0.
 \label{eq:krausiso}
\end{align}
Weyl orthogonality makes its columns orthonormal. Expanding matrix units in the Weyl basis gives
\begin{equation}
 \ket\psi\longmapsto
 \sqrt{1-p}\ket\psi_S\ket{\emptyset}
 +\sqrt p\ket{\Phi_q}_{SA}\ket\psi_B.
 \label{eq:reset}
\end{equation}
The environment has constant dimension $q^2+1$. Only after this conversion may the empty/reset flag be read or copied. Raw photon counting is not this instrument.

\section{Finite stationary processing and the quantum decoder}
\label{sec:digital}

A stored qubit already represents occupations zero and one. Append $b-1$ zero qubits, where $b=\lceil\log_2(Q+1)\rceil$, and use the stored qubit as the least significant bit. This directly embeds $\alpha\ket0+\beta\ket1$ into the binary occupation register without copying, measuring, or transducing the unknown state.

Identical capture commutes with spatial splitting in the linear model. Digitally invert the tree at every Laguerre index using $FN(P-1)$ abstract two-mode gates. Then a computed modal unitary, decomposed into Givens rotations, places the $K$ modes in designated registers using $O(FN^2)$ more. Thus
\begin{equation}
 G_{\rm mode}=O\bigl(F[NP+N^2]\bigr).
\end{equation}
A two-mode gate on total occupation $s\le Q$ is the finite spin-$s/2$ rotation, with tridiagonal entries $\sqrt{(j+1)(s-j)}$. Extend it arbitrarily to larger occupations and padding states. Pair-register dimension is less than $4(Q+1)^2$, so even generic QR synthesis needs only
\begin{equation}
 O\bigl(Q^4\log^2(Q+1)\,\operatorname{polylog}(1/\delta)\bigr)
 \label{eq:paircost}
\end{equation}
bounded few-qubit gates per mode gate. Finite-universal-set precision overhead can use the standard Solovay--Kitaev construction \cite{dawson}. No global Fock or $q^n$ matrix is synthesized. Virtual occupation need not be conserved inside the elementary subcircuit; each completed abstract gate has the correct action on the relevant sector.

In the one-jump comparison, higher occupation bits vanish. Keep the existing low bits and implement \eqref{eq:krausiso}. The same CPTP procedure is defined on every actual output, and discarding other bits contracts comparison error.

For record $\gamma$ with $M$ exports, the preserved theorem represents the instrument by a known unitary history $U_\gamma$ of Hamiltonian segments and mathematical exchanges with fresh maximally mixed qudits. With $d=q^n$ and $b_\gamma=q^M$,
\begin{align}
 \cC_\gamma(\rho)&=\Tr_S U_\gamma(\rho\otimes I_B/b_\gamma)U_\gamma^\dagger,\\
 \cD_\gamma(\sigma)&=\Tr_B U_\gamma^\dagger(I_{S'}/d\otimes\sigma)U_\gamma.
 \label{eq:reverse}
\end{align}
The receiver prepares its own maximally mixed dummy and reverses the known history entirely on receiver registers. This is deterministic CPTP. In normalized operator Hilbert space, $\cD_\gamma\cC_\gamma=J_\gamma^\dagger P_0J_\gamma$, where $J_\gamma(X)=U_\gamma(X\otimes I)U_\gamma^\dagger$ and $P_0$ projects onto source identity. The preserved support-weight estimate gives
\begin{equation}
 \norm{\mathbb E_\gamma\cD_\gamma\cC_\gamma-\id}_\diamond
 \le2\sqrt{n\ee^{-(r_{\rm eff}-\beta)T}}
 \le2\sqrt{n\ee^{-\kappa T/2}}.
 \label{eq:reverseerror}
\end{equation}
Its reference-uniform proof and full-retained-record convention are unchanged.

No fast measurement is assumed. Coherently copy every reset flag into an unused ancilla, run all backward bins in one fixed circuit (inverse Hamiltonian, then flag-controlled export swaps), and trace the copies and consumed systems. For $\sigma=\sum_{\gamma,\gamma'}\ket\gamma\bra{\gamma'}\otimes\sigma_{\gamma,\gamma'}$, tracing the copies removes exactly the off-diagonal record blocks. This equals the measured-record average for \emph{every} acquired state, including finite errors. Another retained copy may provide a classical record.

For $C_H=\sum_{a<b}\norm{[h_a,h_b]}$, first-order product formulas require
\begin{equation}
 O\left(m\left[K+1+\frac{C_H(T+t_0)^2}{\eta_H}\right]+nK\right)
 \label{eq:decodercost}
\end{equation}
fixed-arity gates before synthesis and routing. Small simulated intervals are implemented by bounded physical pulses or finite universal sequences. A bounded-degree layout pays routing; even $O(V)$ swaps per nonlocal gate on $V$ registers remains polynomial.

\section{Finite readout and fixed-response tails}

Every absorber has an adjacent empty storage qubit. At retarded time $t_2=t_1+G_*$, drive
\begin{equation}
 H_{\rm rd}(s)=\Omega(s)(\sigma_a^+\sigma_b^-+\sigma_a^-\sigma_b^+).
\end{equation}
Its full two-qubit norm is $|\Omega(s)|$. Apply a command $\Omega_0$ for $\tau_0=\pi/(2\Omega_0)$ through the same three-pole response. Total pulse area, including the tail, is $\pi/2$. The transfer phase is known and corrected. End finite bookkeeping at $t_3=t_1+\tau_0+2G_*$ and set $A=t_3$. The remaining readout area and residual field coefficient obey
\begin{equation}
 a_{\rm tail}\le\frac{4\Omega_0}{B}\ee^{-BG_*/2},\qquad
 \delta_g\le2g_0\ee^{-BG_*/2}.
\end{equation}

First truncate the actual readout Hamiltonian after $t_3$ as a comparison. Its entire future effect, regardless of future photons, is at most $2D a_{\rm tail}$, where $D=FPN$. Next compare the main capture couplings during readout with zero. A finite-qubit generator bound proves
\begin{equation}
 \norm{\Delta\mathcal L}_\diamond
 \le6F\ell N\sqrt P\,\delta_g+3FPN^2\delta_g^2,
 \qquad\ell=\sqrt r/q.
 \label{eq:tailgenerator}
\end{equation}
Indeed, on a rail write the source coefficient as $S=L/\sqrt P$ and $C=\sum_jc_j$. The dissipator difference contributes at most $4\norm S\norm C+2\norm C^2$. The source--collector and collector--collector cascade Hamiltonians add at most $2\norm S\norm C+N(N-1)\delta_g^2$. Summing rails proves \eqref{eq:tailgenerator}. The common local readout Hamiltonian cancels. Retarded feedforward reordering justifies this bound for the eventual stored-register channel; it does not implement a dense physical cascade Hamiltonian.

Finally the incomplete isolated exchange differs from the intended SWAP by at most $2D a_{\rm tail}$. Thus
\begin{equation}
 \eta_{\rm switch}\le
 (\tau_0+G_*)[6F\ell N\sqrt P\,\delta_g+3FPN^2\delta_g^2]
 +4D a_{\rm tail}.
 \label{eq:switchbudget}
\end{equation}
These are retained source/material channel bounds, with outgoing field traced. They do not assert a bound on every discarded field correlation. Truncation is performed before finite-window photon comparisons; otherwise infinite future emission would make a one-photon-per-rail condition impossible.

The later processor also has fixed poles. For response impulse $h_B(t)=B^3t^2\ee^{-Bt}/2$, a primitive command of area $\theta$ has remaining area at most $|\theta|S_2(Bg)$ after guard $g$, with $S_2(x)=\ee^{-x}(1+x+x^2/2)\le2\ee^{-x/2}$. For $G_{\rm phys}$ routed primitive pulses, truncating tails and charging missing areas gives
\begin{equation}
 \eta_{\rm act}\le4G_{\rm phys}\theta_{\max}S_2(Bg).
\end{equation}
A safe guard is $g=(2/B)\log\max\{\ee,32G_{\rm phys}\theta_{\max}/\eta_{\rm act}\}$. Duhamel's inequality covers noncommuting overlapping tails, and the bound holds at every time after the finite program ends. Gates need no exact OFF or increasing response pole. Wall time is
\begin{equation}
 O\left(G_{\rm phys}\left[\tau_{\max}+B^{-1}\log(G_{\rm phys}/\eta_{\rm act})\right]\right).
 \label{eq:gateclock}
\end{equation}

\section{One explicit noncircular parameter recipe}
\label{sec:recipe}

For $0<\varepsilon\le1$, set $\eta=\varepsilon/32$. Use $\eta$ as the allowance for each separately listed term. Choose
\begin{align}
 T&=\frac2\kappa\log(4n/\eta^2),&
 t_0&=\min\{B^{-1},\eta/(2R_p)\},&t_1&=t_0+T,\\
 K&=\max\{1,\lceil T/\Delta_{\max}\rceil\},&\Delta&=T/K.&&
\end{align}
Here
\begin{equation}
 \Delta_{\max}=\min\left\{\frac1r,\frac\kappa{r^2},
 \frac\eta{2\chi\sqrt{R_pT+(R_pT)^2}},
 \frac{\eta^2}{2nc^2T}\right\}.
\end{equation}
Omit the $\chi$ constraint when $\chi=0$. This pays leading-edge, splitting, one-jump comparison, and reverse-decoder errors. Set $p_Q=(\eta/4)^2$, and define guard-independent constants
\begin{align}
 q_0&=1+(\ee-1)R_p/B+\log(1/p_Q),&z_0&=16q_0^2/\eta^2,\\
 X_0&=g_0^2[t_1+33/(16B)],&P_0&=2+16F(\nu/B)^2/\eta^2,\\
 M_0&=1+K+\frac{256X_0K^2z_0^2}{\alpha^2\Delta^2}
 +\frac{4X_0KK_1^2z_0}{\alpha^4}+2X_0Kz_0,&D_0&=FP_0M_0.
 \label{eq:majorants}
\end{align}
Here $\alpha,K_1$ are from \eqref{eq:rise}. For any trial active length $A$, write $z=1+BA$. Equation~\eqref{eq:cap} gives $Q\le q_0z$. Choose $\zeta=\eta^2/(16Q^2)$, the least integer $N$ satisfying \eqref{eq:modecount}, and the next power of two
\begin{equation}
 P\ge\max\{1,8F\nu^2A^2/\eta^2\}.
 \label{eq:railchoice}
\end{equation}
All roundings are covered by
\begin{equation}
 N\le M_0z^4,\quad P\le P_0z^2,\quad D\le D_0z^6.
\end{equation}
The readout loop closes explicitly. Put
\begin{align}
 a&=1+B(t_1+\tau_0),\\
 C&=\frac{12F\ell g_0M_0\sqrt{P_0}}B
   +\frac{12FP_0M_0^2g_0^2}B+\frac{16D_0\Omega_0}B,\\
 b_*&=4\log\max\{\ee,Ca^{11}/\eta\},\qquad
 y_*=2b_*+2\cdot44^2,\\
 G_*&=y_*/(2B),\qquad A=t_1+\tau_0+2G_*,\qquad z=a+y_*.
 \label{eq:guard}
\end{align}
Compute the actual $Q,\zeta,N,P$ only after this choice. From \eqref{eq:switchbudget}, switching error is at most $Cz^{11}\ee^{-(z-a)/4}$. Since $(a+y)^{11}\le a^{11}(1+y)^{11}$, $\log(1+y)\le\sqrt y$, and Young's inequality gives $y_*\ge b_*+44\log(1+y_*)$, this error is at most $\eta$. The guard is logarithmic in the polynomial coefficients, even when $T$ itself is polynomial. There is no circular hardware-sizing assumption.

All algebraic exponents in this prescription are fixed. For constant local scales and positive margin, source exposure $T=O(\log(n/\varepsilon))$ and the guard remains logarithmic; physical completion can be much longer because of propagation and offline gates. This theorem does not identify source exposure with wall-clock throughput.

\section{Initialization, calibration, noise, and full execution bill}
\label{sec:errors}

\subsection{Input-independent preparation}

For $V$ initially needed qubits, isolated cold ports $\sqrt{\Gamma_p}\sigma^-$ with fixed $\Gamma_p$ reset the registers with channel error at most $2\sqrt{2V}\ee^{-\Gamma_p t_{\rm prep}/2}$. Thus
\begin{equation}
 t_{\rm prep}\ge\Gamma_p^{-1}\log(8V/\eta_{\rm prep}^2)
\end{equation}
suffices. Turn preparation controls off early enough that their residual integrated dissipator also fits the error budget; a fixed-pole guard logarithmic in $V/\eta$ suffices. Signal couplings remain zero during this preparation. Newly allocated processing ancillas can be reset before use; purity, reset hardware, heat removal and time are paid. A thermal reset floor is not removed by waiting.

\subsection{Calibration as a resource}

For coefficient errors bounded by $\delta_g\le g_0$, the retained finite-qubit channel has conservative error
\begin{equation}
 \eta_{\rm wave}\le A[6F\ell N\sqrt P+9FPN^2g_0]\delta_g.
 \label{eq:cal}
\end{equation}
A clock-offset error contributes through $\delta_g\le\delta_{\rm amp}+S_0\delta_t$, with fixed slew $S_0$. Splitter-angle and path-phase errors are bounded by single-particle circuit telescoping, then multiplied by $2\sqrt{R_pA+(R_pA)^2}$. Timing is expressed in calibrated retarded coordinates; this does not assert a uniform trace-norm bound on an arbitrary unlimited-bandwidth field under a time shift. Loss and unmodeled dispersion have separate budgets. Required coefficient precision is inverse-polynomial and program bit length polynomial.

To preserve independent initialization with arbitrary earlier source output, use a one-sided start guard. If the arrival-clock uncertainty is bounded by $\delta_t$, delay the commanded start by that bound. Every actual start then lies in source time $[0,2\delta_t]$, so pre-zero slices see exactly zero coupling. Relative to the nominal schedule this adds at most $2S_0\delta_t$ to the coefficient allowance; readout-clock errors are charged in the local Hamiltonian budget. No vacuum output prehistory is inferred from a vacuum input port.

The theorem can use factory-calibrated components. A conditional offline acquisition of that calibration is explicit: let each fixed-size component have $d_c$ unknown parameters in a known box and $S_c$ bounded-rate probe settings with efficiently computable outcome probabilities $p_s(\theta)$. Assume a specified identifiable branch,
\begin{equation}
 \norm{p(\theta)-p(\theta')}_\infty\ge
 \sigma_{\rm cal}\norm{\theta-\theta'}_\infty,\quad\sigma_{\rm cal}>0,
\end{equation}
with preparation and measurement errors known or jointly identifiable. Also supply a forward bound $\norm{p(\theta)-p(\theta')}_\infty\le L_{\rm cal}\norm{\theta-\theta'}_\infty$, charging $L_{\rm cal}$. Smooth finite-component generators and bounded-duration probes admit such a Duhamel bound. For $M_c$ components, Hoeffding estimates and a fixed-dimensional grid fit achieve accuracy $\delta_\theta$ and total failure probability $\alpha_c$ with
\begin{equation}
 O\left(\sigma_{\rm cal}^{-2}\delta_\theta^{-2}
 \log(2M_cS_c/\alpha_c)\right)
\end{equation}
repetitions per setting. Multiply by components, settings, bounded probe/reset duration and response guards. A grid mesh $O(\sigma_{\rm cal}\delta_\theta/L_{\rm cal})$ needs $O([\operatorname{diam}(\Theta)L_{\rm cal}/(\sigma_{\rm cal}\delta_\theta)]^{d_c})$ points, polynomial for fixed $d_c$ and charged regularity and inverse identifiability. Failure adds at most $2\alpha_c$ to the unconditional channel error. Identifiability, stable drift, and an accurate component model are substantive premises; this is not automatic self-calibration or learning an unknown source Hamiltonian.

\subsection{Uncontrolled noise and storage}

For explicitly added finite-system Hamiltonian errors $H_e$ and inaccessible Lindblad ports $C_e$, a sufficient budget is
\begin{equation}
 \eta_{\rm noise}\le\sum_e\int2\norm{H_e(t)}dt
                  +\sum_e\int2\norm{C_e(t)}^2dt.
 \label{eq:noise}
\end{equation}
It covers full storage lifetimes, including real propagation and all decoding. Channel errors of gates or transport can instead be summed directly. For added hidden source dephasing $\lambda\sum_i\mathcal D[Z_i]$, the conservative extra term over source-time $A$ is $2n\lambda A$. This is a priced perturbation, not a claim that any fixed hidden-noise rate becomes harmless.

The raw capture qubits are uncoded while loading. Later fault-tolerant processing cannot repair arbitrary information already emitted into an inaccessible port. Protected stationary memory may be implemented by a separately specified fault-tolerant architecture, with its encoding, syndrome extraction, ancillas, routing and time charged. No such front-end threshold theorem is assumed here. A fixed uncontrolled loss or drift floor can violate \eqref{eq:noise}; Section~\ref{sec:loss} proves a source-specific consequence.

There is also an explicit attainable imperfect-interface statement. Suppose every original output suffers uniform pure loss of transmissivity $\eta_c$ over the active window $A$. For the source isometry $V_A$, the operator-uniform Poisson law gives
\begin{equation}
 V_A^\dagger\eta_c^{N_{\rm ph}/2}V_A
 =\exp[-R_pA(1-\sqrt{\eta_c})]I.
\end{equation}
This is the overlap of the lossless full source/field purification with the actual beamsplitter-loss purification whose lost input is vacuum. Thus the same deterministic receiver attains the ideal error plus
\begin{equation}
 \eta_{\rm ingress}\le
 2\sqrt{1-\exp[-2R_pA(1-\sqrt{\eta_c})]}
 \le2\sqrt{2R_pA(1-\eta_c)}.
 \label{eq:lossattain}
\end{equation}
No loss event is heralded or discarded. Reducing ingress loss to $O(\varepsilon^2/(R_pA))$ is one sufficient allowance. The converse below gives a different, weaker necessary scaling; no optimal loss threshold is claimed. Internal path loss and capture decay must still satisfy their own complete-channel budgets.

\subsection{Resource ledger and composition}

The principal finite bill is:
\begin{center}\small
\begin{tabular}{p{.29\textwidth}p{.63\textwidth}}
\toprule Resource & Sufficient paid amount\\\midrule
Passive splitters & $F(P-1)$ static components and independent vacuum inputs.\\
Capture and readout & $D=FPN$ two-level absorbers and $D$ adjacent storage qubits.\\
Processing memory & $O(D\log(Q+1)+nK)$ qubits, plus reusable arithmetic/routing workspace.\\
Wire and delay & $O(Fa[P\log P+PN])$ internal wire; maximum internal transit $O(a(P+N)/v)$; actual external delivery added.\\
Mode-processing gates & $O(F[NP+N^2]Q^4\log^2(Q+1)\operatorname{polylog}(F[NP+N^2]Q/\eta))$, before routing.\\
Payload decoder & $O(nK)$ local Kraus/flag circuits and \eqref{eq:decodercost}, plus initial Hamiltonian correction.\\
Classical work & Polynomial Laguerre quadrature, Gram inverse, $N$-dimensional Givens/QR calculations, finite gate synthesis and calibration fitting.\\
Wall time & Preparation and calibration, source window $A$, actual transit, then routed gate clock \eqref{eq:gateclock}. Early stores wait or are protected.\\
Local maxima & Dimension two, capture norm $g_0$, fixed gate norm, fixed reset rate, fixed response poles. Total hardware/power are not fixed.\\\bottomrule
\end{tabular}
\end{center}

The reproducible scalar prescription checks illustrate the severity: at $q=2$, $r=B=g_0=\Omega_0=1$, $\beta=0$, $\chi=1/4$, and $\varepsilon=0.1$, the conservative bound for $n=1$ has $\log_{10}N\simeq98.875$ and $\log_{10}D\simeq112.597$. These are sufficient bounds, not lower bounds or a proposed build. The construction earns a finite operation at bounded local response; it supplies no useful apparatus estimate. The verification file records the exact conventions and also checks $n=4$.

Use mathematical truncation of future readout first, then the finite capture/readout comparison, rail replacement, and cutoff/modal comparison on the actual source. Pass the full-field split and normalized one-jump comparison through the exact-bin extraction and fixed extended decoder. Finally apply the preserved reverse theorem and the early-time inverse. Triangle inequality gives
\begin{align}
 \norm{\cA-\id}_\diamond\le{}&
 \eta_{\rm edge}+\eta_{\rm switch}+\eta_{\rm rail}+\eta_Q+\eta_{\rm modes}
 +\eta_{\rm split}+\eta_{\rm bin}+\eta_{\rm reverse}\nonumber\\
 &+\eta_{\rm prep}+\eta_{\rm numerical}+\eta_{\rm gates}+\eta_{\rm act}
 +\eta_{\rm calibration}+2\alpha_c+\eta_{\rm transport}+\eta_{\rm noise}.
 \label{eq:total}
\end{align}
Assign each displayed term at most $\eta=\varepsilon/32$. The explicit construction leaves slack and yields \eqref{eq:target}. Every comparison is unconditional and reference-uniform. There is no input-dependent favorable-history assumption.

\section{The same-source boundary: collection loss}
\label{sec:loss}

Hardware expansion does not automatically repair a missing incoming quantum subsystem. The following converse uses the actual Weyl ports, not a relabeled cavity-loss or hidden-$Z$ model.

\begin{theorem}[Loss before encoding]
\label{thm:loss}
Take $n$ independent qubit sources with $H=0$ and actual vacuum ports $L_\mu=\sqrt r\,\sigma_\mu/2$, $\mu=X,Y,Z$. Each outgoing port undergoes pure loss of transmissivity $\eta_c$ before any receiver acts. The initial $n$-qubit input is arbitrary and unencoded; the remaining source and lost outputs are inaccessible. For every finite acquisition time and arbitrary collective decoder,
\begin{equation}
 F_{\rm Bell,n}\le[1-c_L(1-\eta_c)^2]^n,\qquad
 c_L=\frac{(1-\ee^{-2})^2}{256}.
 \label{eq:lossmain}
\end{equation}
For any heralded decoder with Bell-input success $s$ and conditional Bell fidelity $F_s$, the same right side bounds $sF_s$.
\end{theorem}

\subsection{A single inaccessible mode retains an initial bit}

The reduced Bloch vector decays as $\ee^{-rt}$. For acquisition time at least $1/r$, take the normalized lost-$Z$ temporal mode
\begin{equation}
 f(t)=\sqrt{\frac{2r}{1-\ee^{-2}}}\ee^{-rt}1_{[0,1/r]}(t),\qquad
 B_f=\int f(t)^*b_{Z,\rm lost,out}(t)dt.
\end{equation}
For initial states $\ket{+Z},\ket{-Z}$ its means are $\pm\mu$, with
\begin{equation}
 \mu^2=\frac{(1-\eta_c)(1-\ee^{-2})}{8}.
\end{equation}
Vacuum input-output regression gives, for all $t,s\ge0$,
\begin{equation}
 \langle b_{\rm lost}^\dagger(t)b_{\rm lost}(s)\rangle
 =\frac{(1-\eta_c)r}{4}\ee^{-r|t-s|}.
\end{equation}
The correlation is input independent. The Schur bound for this kernel gives $N_f=\langle B_f^\dagger B_f\rangle\le(1-\eta_c)/2\le1/2$.

For $Q_f=(B_f+B_f^\dagger)/\sqrt2$, the means are $\pm\sqrt2\mu$, and
\begin{equation}
 \langle Q_f^2\rangle\le N_f+\tfrac12+\sqrt{N_f(N_f+1)}<2.
\end{equation}
Let $p_+,p_-$ be its measurement distributions. Cauchy--Schwarz applied with measure $p_++p_-$ gives
\begin{equation}
 8\mu^2\le4\norm{p_+-p_-}_1.
\end{equation}
Therefore the inaccessible environment distinguishes the initial bit by
\begin{equation}
 D_E:=\norm{\rho_E^+-\rho_E^-}_1
 \ge2\mu^2=\frac{(1-\eta_c)(1-\ee^{-2})}{4}.
 \label{eq:envdist}
\end{equation}
Later dynamics cannot erase this already emitted marginal.

If a recovery has Bell infidelity $e$, purify it. Uhlmann's theorem makes the full recovered pure state within trace distance $2\sqrt e$ of a Bell pair tensor a fixed environment state. Measuring the reference in the $Z$ basis then gives $D_E\le4\sqrt e$. Thus $e\ge D_E^2/16$, proving the one-site version of \eqref{eq:lossmain}. If acquisition time is less than $1/r$, the inaccessible source alone has $Z$-input trace distance $2\ee^{-rT}>2/\ee$, giving the stronger $e>\ee^{-2}/4$. The bound is uniform in time.

\subsection{Collective decoding and heralding are included}

For a normalized state $\rho_{RB}$ with $\dim R=d$, the maximum Bell overlap obtainable by a channel on $B$ satisfies the established recovery/min-entropy SDP \cite{koenig}:
\begin{align}
 dF_{\rm opt}(\rho)&=\min\{\Tr Y:Y\ge0,\ \rho\le I_R\otimes Y\},\\
 &=\max\{\Tr(\rho X):X\ge0,\ \Tr_R X\le I_B\}.
 \label{eq:sdp}
\end{align}
The dual inequality can be completed to equality and then corresponds to a CPTP recovery. Tensor products of primal feasible operators give one product inequality; tensor products of dual feasible operators give the other. Hence
\begin{equation}
 F_{\rm opt}(\rho\otimes\sigma)=F_{\rm opt}(\rho)F_{\rm opt}(\sigma).
\end{equation}
This proves the collective claim; it does not assume a product decoder. Finite-time field Choi states are trace class. Finite-rank truncation and trace-norm continuity of $F_{\rm opt}$ extend the identity to these continuous fields.

A trace-nonincreasing success map can always be completed to a deterministic channel by preparing an arbitrary output on failure. Its unconditional Bell overlap is at least $sF_s$, proving the heralded statement. Success may depend on the input; $s$ is explicitly its value on the Bell input. A bounded-time causal stopping strategy is included by granting the decoder the entire collected field through its latest stopping time.

Consequently, deterministic diamond error is at least
\begin{equation}
 2\{1-[1-c_L(1-\eta_c)^2]^n\}.
\end{equation}
Small diamond error $\varepsilon$ necessarily requires
\begin{equation}
 nc_L(1-\eta_c)^2\le-\log(1-\varepsilon/2).
\end{equation}
Thus $(1-\eta_c)=O(\sqrt{\varepsilon/n})$ is a necessary, not sufficient or optimized, scaling. At fixed loss and fixed nonzero success, conditional Bell fidelity tends to zero. This concerns preserving the entire existing unencoded input, not the capacity for a smaller deliberately encoded logical space.

\section{Fixed aggregate spectral access is a different condition}

The positive theorem fixes local response but grows collective acceptance. A common inaccessible frequency subspace is not removed by adding identical hardware. This too can be shown for the same Weyl source.

At finite acquisition time $T$, zero-extend its emitted field on $[0,T]$ with independent vacuum. Grant the receiver \emph{all} Fourier modes in $|\omega|\le B_f$, even noncausally, and discard their orthogonal complement. This is a generous information-access comparison, not an assertion that an ideal brick-wall filter is physically causal. It applies to architectures whose only accessible input is that projected field; it is not a theorem about every device with a finite response pole.

The spectral projection precedes every receiver operation, including frequency conversion. The specified order is finite-time restriction and vacuum padding, followed by projection. Projection of an ongoing field followed by time restriction is a different operation; the two do not commute.

For initial $\pm Z$, the $Z$-port mean is $\pm m_T(t)$, where $m_T(t)=\sqrt r\,\ee^{-rt}1_{[0,T]}(t)/2$. With Fourier normalization $1/\sqrt{2\pi}$,
\begin{equation}
 |\widetilde m_T(\omega)|^2
 =\frac{r}{8\pi}\frac{|1-\ee^{-(r+i\omega)T}|^2}{r^2+\omega^2}.
\end{equation}
For $T\ge1/r$, the normalized projection of $m_T$ onto the lost frequencies has mean squared at least
\begin{equation}
 \mu_f^2\ge\frac{(1-\ee^{-1})^2}{8}\,h_{B_f},\qquad
 h_{B_f}=1-\frac2\pi\arctan(B_f/r).
\end{equation}
The same correlation-kernel argument gives occupation at most $1/2$ for this normalized lost mode, even though its restriction to $[0,T]$ may have smaller norm. The quadrature and Uhlmann argument therefore proves
\begin{equation}
 e_{\rm Bell}\ge c_Bh_{B_f}^2,\qquad c_B=\frac{(1-\ee^{-1})^4}{256}.
\end{equation}
The inaccessible-source argument covers $T<1/r$. All fields used here are finite-time states; no infinite-photon Fock-state assumption enters. Tensorization and completion give
\begin{equation}
 F_{\rm Bell,n},\ sF_s\ \le(1-c_Bh_{B_f}^2)^n.
\end{equation}
For small requested diamond error,
\begin{equation}
 B_f=\Omega\!\left(r\sqrt{n/\varepsilon}\right)
\end{equation}
is necessary in this spectral-access model. An exact finite bound is
\begin{equation}
 B_f\ge r\cot\!\left[\frac\pi2
 \sqrt{\frac{-\log(1-\varepsilon/2)}{nc_B}}\right]
\end{equation}
when the square root is less than one. This is a necessary aggregate-access bound, not a local-linewidth lower bound.

\subsection{A sharp ordinary amplitude-emission comparator}

For a source emitting $\ket0\mapsto\ket{\mathrm{vac}}$, $\ket1\mapsto\ket{1_f}$, retaining squared mode fraction $\eta$ gives amplitude damping. Its exact optimum is
\begin{equation}
 F_{\rm Bell,n}^{\rm opt}=\left[\frac{(1+\sqrt\eta)^2}{4}\right]^n.
\end{equation}
For one copy the positive witness $Y=\tfrac12\operatorname{diag}(1+\sqrt\eta,\eta+\sqrt\eta)$ obeys $\rho_{RB}\le I_R\otimes Y$; identity recovery attains $\Tr Y/2$. Tensoring proves the collective statement. For the ordinary exponential mode $f(t)=\sqrt\gamma\ee^{-\gamma t/2}1_{t\ge0}$,
\begin{equation}
 \eta_{B_f}=\frac2\pi\arctan(2B_f/\gamma),
 \qquad B_f=\Omega(n\gamma/\varepsilon)
\end{equation}
is necessary for full-input diamond accuracy. This sharper formula belongs to the amplitude-emission model. It is not substituted for the Weyl-source loss law.

\section{Preserved extensions, stopping, and boundaries}

For Hamiltonians diagonal in one known fixed product basis, the previous exact construction measures each reset purifier $A$ in that basis, retains the first export per site, and tracks the later replacement labels. Conditioned on the record, the retained initial payload is multiplied by a known phase
\begin{equation}
 \exp\!\left[-i\int\sum_X h_X(t;z_X(t;x))dt\right].
\end{equation}
One final fixed-arity diagonal gate per Hamiltonian term removes it. For every $r>0$ the ideal uncompleted-input error is at most $2n\ee^{-rT}$. In the present bin comparison use $r_{\rm eff}$, or ensure $r_{\rm eff}\ge r/2$ and take $T\ge(2/r)\log(2n/\eta)$. All remaining capture and execution budgets are as above; strength still enters $\chi$, precision and gates. This preserves the exact diagonal mechanism and supplies the bounded-component interface. General noncommuting Hamiltonians with $r\le\beta$ remain unresolved.

The preceding observable stopping certificate is unchanged: a positive $n$-variable process is propagated between resets and zeroed at a reset site, bounding each ideal record's recovery error. Its stronger sufficient rate and conditional-error convention remain intact. This batch modal receiver discovers flags after capture; a retrospectively certified prefix does not automatically shorten the already completed physical loading time. No online stopping advantage is silently imported.

The previous interaction-order theorem retains its actual summed noncommuting-term budget and full-retained-instrument scope. Its weak-control proof does not decide the strong-control regime. The unencoded hidden-$Z$ obstruction remains distinct from accessible coherent Weyl emission. Preparing entanglement first and teleporting a later payload is a different operation and is not contradicted by the present loss bounds.

Two mixed-output distinctions also remain. N1733 requires an exact extreme channel, an uncorrelated finite mixed bath, and access to the whole relevant bath; its positive whole-bath decoder is not automatically causal \cite{astra1733}. The separate preserved expander companion concerns a time-programmed source with maximally mixed collision bath and unavailable purifiers. Its whole-output storage obstruction is not a theorem about an arbitrary stationary thermal Lindbladian. This report acquires the initial payload and its reference correlations, not every emitted bath correlation.

\section{Verification and provenance}

The proof was divided among independent constructive branches: continuous input/output and finite-response controls; qubit compilation and loss converses; and stationary quantum processing. Cross-audits identified and corrected two material bookkeeping errors before integration: graph depth was replaced by actual fixed-density wire delay, and the photon-cap comparison was placed on the actual source before normalizing good histories. A separate finite-qubit generator audit closed both capture and readout tails without a circular sizing assumption.

The supplied reproducible numerical checks test exact Laguerre/cascade identities, the finite vacuum Gram identity, local Kraus conversion, actuator bounds, the amplitude-damping witness, and the explicit finite parameter recipe. They are diagnostics, not proofs of the asymptotic theorem or interval-certified hardware performance. The analytical constructions above and the full companion notes carry the claims. No historical-priority or external peer-review claim is made.

The preceding many-body package is copied without alteration. Its SHA-256 is
\begin{center}\ttfamily\footnotesize
c7bc392421457f3c807516bdd06b9458954e1cb621f62f8cfbb826d69507cb55.
\end{center}
It in turn preserves the interaction-order theorem, meter constructions, digital causal-reset opening, mixed-output companion, earlier handoffs, and verification. The new package includes the current N1732/N1733 governing cards and the retrieved complete visible source return. Their supplied-source status remains distinct from the constructive work in this report.

\begin{thebibliography}{9}
\bibitem{yamamoto} N. Yamamoto and M. R. James, ``Zero-dynamics principle for perfect quantum memory in linear networks,'' \emph{New Journal of Physics} \textbf{16}, 073032 (2014). \href{https://doi.org/10.1088/1367-2630/16/7/073032}{doi:10.1088/1367-2630/16/7/073032}; \href{https://arxiv.org/abs/1403.1698}{arXiv:1403.1698}.
\bibitem{nurdin} H. I. Nurdin and J. E. Gough, ``Modular Quantum Memories Using Passive Linear Optics and Coherent Feedback,'' \emph{Quantum Information and Computation} \textbf{15}, 1017--1040 (2015). \href{https://arxiv.org/abs/1409.7473}{arXiv:1409.7473}.
\bibitem{reck} M. Reck, A. Zeilinger, H. J. Bernstein and P. Bertani, ``Experimental realization of any discrete unitary operator,'' \emph{Physical Review Letters} \textbf{73}, 58--61 (1994). \href{https://doi.org/10.1103/PhysRevLett.73.58}{doi:10.1103/PhysRevLett.73.58}.
\bibitem{dawson} C. M. Dawson and M. A. Nielsen, ``The Solovay--Kitaev algorithm,'' \emph{Quantum Information and Computation} \textbf{6}, 81--95 (2006). \href{https://arxiv.org/abs/quant-ph/0505030}{arXiv:quant-ph/0505030}.
\bibitem{koenig} R. K\"onig, R. Renner and C. Schaffner, ``The operational meaning of min- and max-entropy,'' \emph{IEEE Transactions on Information Theory} \textbf{55}, 4337--4347 (2009). \href{https://doi.org/10.1109/TIT.2009.2025545}{doi:10.1109/TIT.2009.2025545}; \href{https://arxiv.org/abs/0807.1338}{arXiv:0807.1338}.
\bibitem{astra1732} Astra Research, governing packet N1732, ``Collect a known Markov source's actual pure ports with a finite causal receiver,'' current content version \texttt{86bbc29b178a...}. \href{https://github.com/harjassand/astra-research-knowledge/blob/main/knowledge/current/packets/N1732.json}{Public governing card}. Full visible source return and unchanged card included with provenance.
\bibitem{astra1733} Astra Research, governing packet N1733, ``Recover an extreme channel complement from the whole finite mixed bath,'' current content version \texttt{2144ec8db88c...}. \href{https://github.com/harjassand/astra-research-knowledge/blob/main/knowledge/current/packets/N1733.json}{Public governing card}. Whole-bath, exact-extremality and noncausal qualifications retained.
\end{thebibliography}
\end{document}
